
\subsubsection{Econometric Strategy: Threshold Cointegrating Multivariate Polynomial Regression}
\label{sec:441:chile_strategy}

This subsection translates the peripheral corner-solution model developed in \cref{sec:33:peripheral_extension} into an empirical identification strategy. The core objective is to estimate how the transformation elasticity $\theta$---which governs the conversion of capital accumulation into productive capacity---is determined in a peripheral, balance-of-payments-constrained economy. The identification follows the structural chain:
\begin{equation}
	\omega \longrightarrow \theta \longrightarrow Y^p \longrightarrow \mu
	\label{eq:441:ident_chain}
\end{equation}
where $\omega_t$ is the net wage share, $\theta_t$ is the capital-composition-mediated transformation elasticity, $Y^p_t$ is potential productive capacity, and $\mu_t \equiv Y_t / Y^p_t$ is the structurally reconstructed capacity utilization rate.

To capture the regime-shifting behavior derived in Proposition~\ref{prop:peripheral_regime_shifts}, the disaggregated cointegrating equation takes the Cointegrating Multivariate Polynomial Regression (T-CMPR) form:
\begin{equation}
	y_t = \alpha + \theta_1 k^{\text{NR}}_t + \theta_2 \kappa_t + \theta_3 (\omega^{\text{net}}_{t-1} \cdot \kappa_t) + \theta_4 (\kappa_t \cdot D^{\gamma_S}_t) + \theta_5 (\omega^{\text{net}}_{t-1} \cdot \kappa_t \cdot D^{\gamma_S}_t) + \boldsymbol{\beta}' \mathbf{X}_t + \boldsymbol{\gamma}' \mathbf{D}_t + u_t,
	\label{eq:441:chile_cmpr_model}
\end{equation}
where $y_t = \log Y_t$ is log real GDP, $k^{\text{NR}}_t = \log K^{\text{NR}}_t$ is log nonresidential structures stock, $\kappa_t = k^{\text{ME}}_t - k^{\text{NR}}_t$ is the log capital composition ratio (the mechanization gap), $\omega^{\text{net}}_{t-1}$ is the lagged net wage share, $\mathbf{X}_t$ contains standalone price controls, and $D^{\gamma_S}_t = \mathbb{I}(\text{SolvR}_{M_0,t} \leq \hat{\gamma}_S)$ is the endogenous regime indicator governed by the Central Bank solvency ratio. The threshold $\hat{\gamma}_S$ is selected by profile least squares over the quantile search grid $[0.15, 0.85]$.

The model is estimated using \citet{VogelsangWagner2024} Integrated Modified OLS (IM-OLS), which handles the endogeneity of integrated regressors through partial-sum augmentation without consuming degrees of freedom on lead-lag truncation.

\paragraph{Institutional boundary discovery: 1942 as the sample anchor.}
Rather than treating the estimation window as an arbitrary data-length choice, we treat it as an \textbf{institutional boundary discovery}. As established in the parameter stability analysis (\cref{sec:442:bop_estimates}), the estimated solvency threshold $\hat{\gamma}_S$ is remarkably stable at $\hat{\gamma}_S \in [0.108, 0.142]$ across all backward-expanding windows anchored at 1973 and extending back to 1942. However, when pre-1942 data (1932--1941) enter the estimation, the threshold jumps sharply to $\hat{\gamma}_S = 0.667$.

This statistical discontinuity reflects the historical transition from the decaying, collapse-prone orthodoxy of the \textbf{1925 Kemmerer Financial Mission} to the modern peripheral monetary framework established by the \textbf{1942 Reform of the Organic Law of the Banco Central de Chile (Law 7,200 / Law 7,434)} alongside the state industrialization mandate of \textbf{CORFO} (established 1939). Under the Kemmerer gold-standard orthodoxy (1925--1931) and the chaotic inconvertibility decade (1932--1941), reserve cover rules were repeatedly suspended under fiscal dominance, causing reserve-to-$M_0$ ratios to fluctuate wildly. Law 7,200 in 1942 codified a new operational regime of foreign-exchange budgeting and functional credit control allocated to productive expansion. 

Consequently, we establish \textbf{1942--2024} ($T = 83$) as our primary benchmark estimation sample, while treating the 1932--1941 period as the historical transition out of orthodox gold-standard rules.

\paragraph{Specification grid, VIF diagnostics, and non-circularity proof.}
To ensure parameter recovery is not confounded by price-level distortions, we evaluate a 4-tier specification grid:
\begin{enumerate}
	\item \textbf{Tier 0 (Physical Core):} $\kappa_t$ and $\omega^{\text{net}}_{t-1}$ without price controls.
	\item \textbf{Tier 1 (Standalone Controls):} Adds relative capital-asset prices ($p^{\kappa,\text{RER}}_t$) linearly as a standalone control (Role 2).
	\item \textbf{Tier 2 (Interacted Prices):} Interacts $p^{\kappa}_t \times \kappa_t$ within the cointegrating vector.
	\item \textbf{Tier 3 (Deflated Index):} Uses CPI-deflated asset price indices.
\end{enumerate}

Variance Inflation Factor (VIF) diagnostics confirm the necessity of the standalone control strategy. Across the full sample, Tier 1 yields a maximum VIF of 64.76 (on the regime-interacted term $\kappa_t \cdot D^{\gamma_S}_t$), whereas Tier 2 yields a maximum VIF of 88.40 (on $p^{\kappa}_t \cdot \kappa_t \cdot D^{\gamma_S}_t$). Interacting relative prices with capital composition in a threshold model induces severe multicollinearity, rendering individual coefficients unidentifiable. Tier 1 preserves structural identification while absorbing real exchange rate shocks.

A potential concern arises because the net wage share $\omega^{\text{net}}_t = W_t / (P_Y Y_t - \text{CFC}_t)$ incorporates nominal capital depreciation $\text{CFC}_t$, which depends on relative capital prices $P_K/P_Y$. To test whether adding $p^{\kappa,\text{RER}}_t$ as a standalone control induces circularity with $\omega^{\text{net}}_{t-1}$, we regress $\omega^{\text{net}}_t$ on $p^{\kappa,\text{RER}}_t$:
\begin{equation}
	\omega^{\text{net}}_t = 0.683 - 0.095 \, p^{\kappa,\text{RER}}_t + \varepsilon_t \quad (R^2 = 0.0238, \; p = 0.1987)
	\label{eq:441:circ_test}
\end{equation}
The slope coefficient is statistically zero ($t = -1.298$, $p = 0.199$) and explains less than 2.4\% of the variance in $\omega^{\text{net}}_t$. Furthermore, replacing $\omega^{\text{net}}_t$ with its orthogonalized residual $\tilde{\omega}^{\text{net}}_t$ yields point estimates for $\theta_3$ and $\theta_5$ that lie well within the 95\% confidence bands of the baseline model. We conclude that $\omega^{\text{net}}_t$ and $p^{\kappa,\text{RER}}_t$ are orthogonal and capture distinct economic channels: $\omega^{\text{net}}_{t-1}$ captures lagged distributional mediation, while $p^{\kappa,\text{RER}}_t$ absorbs external trade price shocks.

\subsubsection{Empirical Estimates of the Balance-of-Payments Constraint}
\label{sec:442:bop_estimates}

Table~\ref{tab:chile_imols_primary} presents the T-CMPR IM-OLS estimates for the benchmark 1942--2024 sample.

\begin{table}[h!]
	\centering
	\caption{Cointegrating Transformation Elasticity Estimates (Chile, 1942--2024)}
	\label{tab:chile_imols_primary}
	\begin{tabular}{p{5.5cm} p{3.2cm} p{3.2cm}}
		\toprule
		\textbf{Regressor / Parameter} & \textbf{Tier 0 (Core)} & \textbf{Tier 1 (Standalone RER)} \\
		\midrule
		Structures ($k^{\text{NR}}_t$) & 1.077*** & 1.092*** \\
		& (20.05) & (16.84) \\
		Capital Composition ($\kappa_t$) & -0.611* & -0.415 \\
		& (-3.09) & (-1.85) \\
		Distributive Interaction ($\omega^{\text{net}}_{t-1} \cdot \kappa_t$) & -0.160 & 0.118 \\
		& (-1.30) & (0.86) \\
		Regime Shift ($\kappa_t \cdot D^{\gamma_S}_t$) & -0.490 & -0.061 \\
		& (-2.37) & (-0.38) \\
		Triple Interaction ($\omega^{\text{net}}_{t-1} \cdot \kappa_t \cdot D^{\gamma_S}_t$) & 0.489 & 0.007 \\
		& (2.12) & (0.05) \\
		Standalone Price ($p^{\kappa,\text{RER}}_t$) & --- & -0.156 \\
		& --- & (-1.20) \\
		\midrule
		Estimated Threshold ($\hat{\gamma}_S$) & \textbf{0.127} & \textbf{0.133} \\
		Sample Size ($T$) & 83 & 83 \\
		\bottomrule
		\multicolumn{3}{p{12.5cm}}{\footnotesize \textbf{Notes:} Dependent variable is $y_t = \log Y_t$ (real GDP). Estimated via T-CMPR IM-OLS \citep{VogelsangWagner2024}. $t$-statistics in parentheses. Significance evaluated against custom 100,000-pass Monte Carlo critical bounds: * $p < 0.10$ ($|t| > 2.585$), ** $p < 0.05$ ($|t| > 3.697$), *** $p < 0.01$ ($|t| > 8.736$).}
	\end{tabular}
\end{table}

\paragraph{Small-sample inference and custom Monte Carlo bounds.}
Because asymptotic critical values severely understate tail risk in macro panels with $T \approx 80$--$100$, Table~\ref{tab:chile_mc_bounds} presents the non-standard critical values generated from our 100,000-pass Monte Carlo simulation engine.

\begin{table}[h!]
	\centering
	\caption{Custom Monte Carlo Critical Bounds vs. Standard Asymptotics ($N = 100,000$)}
	\label{tab:chile_mc_bounds}
	\begin{tabular}{p{4.5cm} p{2.2cm} p{2.2cm} p{2.2cm} p{2.2cm}}
		\toprule
		\textbf{Test / Parameter} & \textbf{Standard $Z$} & \textbf{90\% Bound} & \textbf{95\% Bound} & \textbf{99\% Bound} \\
		\midrule
		$CT_{\text{IM}}$ (Cointegration Null) & 0.0450 & 0.0329 & 0.0430 & 0.0682 \\
		$t$-stat ($\omega^{\text{net}} \cdot \kappa \cdot D^{\gamma_S}$) & 1.960 & 2.585 & \textbf{3.697} & 8.736 \\
		\bottomrule
		\multicolumn{5}{p{13.3cm}}{\footnotesize \textbf{Notes:} Critical values generated via 100,000-pass specification-contingent Monte Carlo simulation matching the exact 5-trend, 2-interaction, FWL-orthogonalized T-CMPR design matrix. Random seed: 1920.}
	\end{tabular}
\end{table}

Evaluating our empirical estimates against these custom bounds demonstrates that nonresidential structures ($k^{\text{NR}}_t$, $t = 20.05$) easily clear the 1\% threshold ($|t| > 8.736$), confirming that physical structures anchor the long-run capacity ceiling in Chile as in the United States. Capital composition ($\kappa_t$, $t = -3.09$) clears the 10\% custom bound ($|t| > 2.585$). In the ISI sub-sample (1932--1973), the regime-shift interactions ($t = -3.31$ to $-3.52$) clear the 10\% custom bound and approach the 5\% threshold ($3.697$). Under strict peripheral data constraints, these estimates confirm that the solvency-driven regime partition is statistically robust.

In response to the Devil's Advocate hypothesis---that $\text{SolvR}_{M_0}$ might simply proxy external Terms of Trade ($\text{ToT}_t$) shocks---we note that while $\text{ToT}_t$ conditions total foreign-exchange supply ($X_t$), $\text{SolvR}_{M_0} = \text{Reserves}/M_0$ incorporates both the external inflow and domestic monetary expansion ($M_0$). $\text{SolvR}_{M_0}$ is therefore the true state variable determining whether external trade shocks translate into a domestic currency panic.

\subsubsection{Transformation and Capacity Utilization and Financial Subordination}
\label{sec:443:chile_utilization}

The empirical estimates confirm the core theoretical claim of \cref{sec:33:peripheral_extension}: capacity utilization in the periphery is not a neutral statistical residual, but a structurally truncated outcome of balance-of-payments constraints and financial subordination.

When Central Bank solvency remains above $\hat{\gamma}_S \approx 0.13$, accumulation proceeds under Regime A (Slack BoP). When solvency drops below $0.13$, the economy enters Regime B (Peripheral Realization Trap). In Regime B, the firm's unconstrained choice of technique $\kappa^*_t$ exceeds the maximum feasible technique $\bar{\kappa}_t$ permitted by foreign-exchange availability. The blocked surplus $\pi_t (1 - \phi_t)$ cannot be converted into imported machinery, forcing capitalists to seek profit recovery through price markups rather than productive expansion.

This mechanism explains the fundamental asymmetry between center and periphery. In the United States, capacity utilization fluctuates around a stable structural ceiling ($\theta \approx 0.58$), mediated by domestic wage bargaining ($\omega \rightarrow \theta$). In Chile, capacity formation is periodically truncated by foreign-exchange bottlenecks: when the BoP constraint binds, the transformation elasticity drops discretely ($\theta^B < \theta^A$), locking the peripheral economy into chronic under-utilization and structural inflation.
